% content/08_conclusion.tex
\phantomsection%
\section*{Заключение}
\addcontentsline{toc}{section}{Заключение}

В первой лабораторной работе было создано представление для подсчёта заселений и платежей каждого студента. Запрос для сравнения числа платежей вернул 4462 строки. Запрос с подсчётом отзывов о проживании вернул 5000 строк.

Во второй работе была создана сводная таблица и реализовано семь триггеров для её автоматического обновления. Во всех семи проверках состояние сводной таблицы после изменения исходных данных совпало с ожидаемым.

В третьей работе были настроены две роли с правом входа. Роль \texttt{dorm\_operator} позволяла изменять записи о студентах, заселениях и платежах, а роль \texttt{dorm\_viewer} --- просматривать обновлённые сводные данные. Запросы, выходившие за пределы выданных полномочий, завершались отказом в доступе.

В четвёртой работе в двух параллельных транзакциях на уровне изоляции \texttt{READ COMMITTED} было воспроизведено неповторяющееся чтение. После фиксации изменения одной транзакцией повторное чтение в другой показало, что сумма платежа изменилась с 693,88 до 1000,00.

В пятой работе были реализованы функция \texttt{format\_full\_name} для вывода фамилии и инициалов студента и процедура \texttt{register\_settlement} для регистрации заселения и связанного с ним платежа. Работа процедуры была проверена в трёх сценариях с разным набором уже имеющихся записей.

Для работы с PostgreSQL применялся консольный клиент \texttt{psql}. При разработке использовались языки SQL и PL/pgSQL\@. Схема базы данных подготовлена в draw.io, а блок-схема процедуры --- в Graphviz. Отчёт собран средствами LaTeX на компьютере с macOS и процессором архитектуры ARM64.
